\documentclass[a4paper]{article}
\usepackage[pdftex]{graphicx}
\usepackage[utf8]{inputenc}
\usepackage{enumerate}
\usepackage{siunitx}
\sisetup{locale=DE}
\usepackage{href-ul}
\hypersetup{
	colorlinks=true,
	linkcolor=blue,
	urlcolor=blue}
\usepackage{icomma}
\usepackage{amssymb}
\hyphenation{Tem-pe-ra-tur-ska-la}
\usepackage{geometry}
\geometry{a4paper, top=15mm, left=15mm, right=15mm, bottom=15mm,
	headsep=10mm, footskip=12mm}

\begin{document}
	%Title
	{\bf School assignment from physics }\\
	\begin{enumerate}[1.]
		%Task 1
		{\bf \item Explain what is meant by the following terms:}
		
		a) Heat \hspace{1 cm} b) Internal energy \hspace{1 cm} c) Anomaly of water
		
		%Exercise 2
		{\bf \item Name three different effects that changing the internal energy can have.}
		
		%Task 3
		{\bf \item Why is it possible to set a temperature scale using the Anders Celsius method?}
		
		%Task 4
		{\bf \item Describe the process of heat conduction in the particle model.}
		
		%Task 5
		{\bf \item Describe in detail the development of the heat flow of air in a room that is heated with the help of a radiator.}
		
		%Task 6
		{\bf \item Tea in a thermos stays hot for a long time.} \\
		Describe the physical measures by which this is achieved and briefly explain why this is possible in each case.
		
		%Task 7
		\begin{minipage}{0.65\textwidth}
			{\bf \item In an experiment, the connection between the amount of heat supplied to a body $W_{th}$} and the change in temperature of the body is examined. To do this, friction work $W_R$ is performed on a metal cylinder with diameter $d=\SI{5,6}{\centi\meter}$ by rotating the cylinder. The string that is wrapped around the cylinder is tensioned by a mass with the weight $F_G= \SI{25}{\newton}$. n is the number of revolutions.
		\end{minipage}
		\hfill
		\begin{minipage}{0.3\textwidth}
			\includegraphics[width=4cm]{rolle112.pdf}
		\end{minipage}
		
		\begin{enumerate}[a)]
			\item Calculate the amount of heat supplied to the cylinder during one revolution. \\
			
			The following measurement table results:
			
			\renewcommand{\arraystretch}{2}
			\begin{tabular}{ccccc}
				$\Delta \vartheta ~\textnormal{in}\SI{}{\celsius}$ & 0.51 & 0.98 & 1.5 & 2.3 \\
				\hline
				n & 100 & 200 & 300 & 450 \\
				\hline
				$W_{th} \textnormal{in}~\SI{}{\kilo\joule}$ & & & &
			\end{tabular}
			
			\item Complete the table with $W_{th} ~\textnormal{in}\SI{}{\kilo\joule}$ and evaluate it graphically.\\
			$[\textnormal{y-axis}~= W_{th}; \SI{1}{\centi\meter}~ \hat{=}~\SI{0.25}{\kilo\joule};~
			\textnormal{x-axis}~=\Delta \vartheta ;~ \SI{1}{\centi \meter}~ \hat{=}~\SI{0.25}{\celsius}] $
			What is the connection?\\
			
			\item How can you recognize the connection mathematically? What is the mass of the cylinder if it is made of brass with the specific heat capacity
			$c=\SI{0.384}{\kilo\joule\over \kilogram\cdot\kelvin}$ is?
		\end{enumerate}
		
	\end{enumerate}
	
	\fbox{
		\begin{minipage}{0.5\textwidth}
			For the solution please \href{https://www.okuyakl.de/phys/p9thcL112/le112.pdf}{click here} or scan the QR code.\\
			Additional worksheets are available at
			
			\href{http://www.okuyakl.com}{www.okuyakl.com}
		\end{minipage}
		\hfill
		\begin{minipage}{0.4\textwidth}
			\includegraphics[width=1.5 cm]{../../viecher/zwe03}
			\includegraphics[width=3 cm]{qre112}
			\includegraphics[width=2 cm]{../../viecher/afanticon1}
			
	\end{minipage}}
	
\end{document}